\documentclass{standalone}
\usepackage{standalone}

\begin{document}
\chapter{Conclusion}

In this project I developed Sarabo, a platform that turns informal electronics repair, which today depends almost entirely on trust, into a tracked service with clear prices and a full record of every step. Customers describe the problem in their own language and compare ranked offers from qualified local technicians. They pay an inspection fee that is later taken off the final price, approve an itemised quote, follow every step of the repair, and are protected by a 30-day warranty. Technicians get a steady source of jobs that match their skills and a wallet they can see into. Admins get a single view of everything that needs their attention.

The main engineering result is a workflow in which the server has the final word. An action only succeeds if the caller has the right role, the request is in the right state, and every amount comes from stored data. Transactions, idempotent payments and refunds, and balances that are calculated instead of stored keep money and state consistent even when actions are retried or happen at the same time. Every waiting state now has a deadline, an automatic step or an admin alert, so no repair can get stuck. Generative AI is used only as an assistant, and its output is always validated and reviewed by a person.

Since the earlier report, Sarabo has grown from a workflow in which admins assigned every job into a live two-sided marketplace. Over the same period its automated server checks grew from 1,890 to \servertestcount{}, all passing, alongside \clienttestcount{} browser tests. The remaining limitations (Chapter~\ref{ch:limitations}) show what still lies between a validated academic system and a live service. The most important of these are a local payment gateway, automatic payouts and use by real customers.

\end{document}
